\documentclass[exercice,a4]{thedocument}%
\usepackage{texmaths-tikz}%
\usepackage{pstricks}%
\usepackage{pst-bezier}%
\usepackage{pstricks-add}%
\usepackage{pst-eucl}%
\usepackage{pst-func}%
\usepackage{pst-tree}%
\usepackage{pst-circ}%
\usepackage{tikz}%
\usepackage{pgfplots}%
\usepackage{tkz-tab}%
\usepackage{tikzpagenodes}%
\usepackage{tabularray}%
\usepackage{tabularx}%
\usepackage{multirow}%
\usepackage{multicol}%
\usepackage{scratch3}%
\usepackage{enumitem}%
\usepackage[np]{numprint}%
\usepackage{pas-tableur}%
\usepackage{listings}%
\usepackage{graphicx}%
\usepackage{adjustbox}%
\usepackage{needspace}%
\usepackage{xcolor}%
\usepackage{eurosym}%
\usepackage{diagbox}%
\usepackage{colortbl}%
\usepackage{multido}%
\usepackage{icomma}%
\usepackage[french]{babel}%
\graphicspath{{../brevet/2026/}}%
\startdatakeys
    \source{26GENMATMEAG1}
    \BO{2026}
    \cycle{4}
    \competencesDuSocle{CA,RA,CO}
    \objectifsApprentissage{C4A1aa,C4A2ad,C4B2cb,C4D2aa,C4D3ac}
    \motsCles{facteurs premiers,PGCD}
\enddatakeys
\definecolor{gris}{gray}{0.5}%
\providecommand{\wideparen}[1]{\overset{\frown}{#1}}%
\begin{document}%
\providecommand{\degres}{^\circ}%
On souhaite réaliser des sachets de bonbons.

Pour cela, on dispose de 112 bonbons à la fraise et de 140 bonbons au caramel.

\begin{itemize}[left=5mm]
	\item Tous les sachets contiennent le même nombre de bonbons à la fraise.

	\item Tous les sachets contiennent le même nombre de bonbons au caramel.

	\item Tous les bonbons à la fraise et au caramel sont utilisés.
\end{itemize}

\begin{enumerate}
	\item Peut-on constituer 16 sachets?

	\item La décomposition en facteurs premiers de 112 est $2^{4} \times 7$.

	Quelle est la décomposition en facteurs premiers de 140 ?

	\item Quel nombre maximal de sachets pourra-t-on constituer ?

	Quelle sera alors la composition de chaque sachet?
\end{enumerate}

% Fin du document
\medskip
\showthesource
\end{document}
